\subsection{无任意小有限子群的群的情形}

定理 2. --- 设 G 为满足下列条件的局部紧群：

\begin{enumerate}
\def\labelenumi{(\Alph{enumi})}
\setcounter{enumi}{11}
\tightlist
\item
  存在 \(e\) 在 \(G\) 中的一个邻域，它不含 \(G\) 的任何不约化为 \(e\) 的有限子群。
\end{enumerate}

于是下列性质成立：

\begin{enumerate}
\def\labelenumi{(\roman{enumi})}
\item
  G 的离散子群所成的集合 D 在 \(\Sigma\) 中是局部闭的（这等价于说它是局部紧的）。
\item
  为使 \(\mathbf{A}\) （\(\mathbf{D}\) 的闭子集）是紧的，必须且只需存在 \(\mathbf{U}\) （\(e\) 在 \(\mathbf{G}\) 中的一个邻域），使得 \(\mathbf{H} \cap \mathbf{U} = \{e\}\) 对每个子群 \(\mathbf{H} \in \mathbf{A}\) 都成立。
\item
  若此外 G 还由 e 的一个紧邻域生成，则由 G 中使 G/H 为紧的离散子群 H 所成的集合 \(D_{c}\) 在 \(\Sigma\) 中是局部闭的（从而是局部紧的）。
\end{enumerate}

我们有 \(\mathbf{D}_c = \mathbf{D} \cap \Sigma_c^0\) ，因此 (iii) 是 (i) 以及 No.~3 的 Th. 1 (ii) 的推论。

采用 No.~3 命题 7 的记号，为证 (i) 与 (ii)，只需证明：

命题 8. --- N 到 D 上的典范双射是一个同胚。

现在，若 \(\Gamma_d\) 是 G 的离散子群上的 Haar 测度所成的集合，则 D 典范地同胚于 \(\Gamma_d\) 中比 \(>0\) 的位似群的轨道空间 (GT, I, §5, No.~2, Prop. 4)。因此只需证明典范映射 \(\alpha \mapsto (\alpha(\{e\}), \alpha/\alpha(\{e\}))\) 是 \(\Gamma_d\) 到 \(\mathbf{R}_+^* \times \mathbf{N}\) 上的一个同胚，而这将由下面的引理得出：

引理 4. --- 若群 G 满足条件 (L)，则映射 \(\alpha \mapsto \alpha (\{e\})\) （\(\Gamma_d\) 到 \(\mathbf{R}_+^*\) 中的）是淡连续的。

考察一个测度 \(\alpha \in \Gamma_d\) ；设 \(V_0\) 是 \(e\) 在 \(G\) 中的一个相对紧开邻域，使得 \(H_\alpha \cap V_0 = \{e\}\) ，并且不含 \(G\) 的任何包含在 \(V_0\) 中的不约化为 \(e\) 的有限子群。设 \(V\) 是 \(e\) 的一个对称紧邻域，使得 \(V^3 \subset V_0\) ，又设 \(U\) 是 \(e\) 的一个对称邻域，使得 \(U^2 \subset V\) 。设 \(\varphi\) （相应地 \(\psi\) ）是 \(\mathcal{K}_+(G)\) 中的一个函数，取值于 \([0,1]\) ，在 \(V^3\) 上（相应地在点 \(e\) 处）等于 1，且支集包含在 \(V_0\) 中（相应地包含在 \(U\) 中）。\(\beta \in \Gamma_d\) 中使得 \(|\beta(\varphi) - \alpha(\varphi)| \leqslant \varepsilon\) 与 \(|\beta(\psi) - \alpha(\psi)| \leqslant \varepsilon\) 的测度所成的集合，是 \(W\) ，即 \(\alpha\) 的一个邻域。我们要证明：只要 \(\varepsilon\) 取得充分小，\(H_\beta \cap V = \{e\}\) 对每个 \(\beta \in W\) 都成立；由此将推出 \(\beta(\psi) = \beta(\{e\})\) ，从而 \(|\beta(\{e\}) - \alpha(\{e\})| \leqslant \varepsilon\) ，这就证明了引理。

只需证明：对 \(\beta \in \mathbf{W}\) ，

\[
(\mathrm{V} ^ {2} - \mathrm{V}) \cap \mathrm{H} _ {\beta} = \varnothing .\tag{8}
\]

因为，假设这一点已经确立：那么对 \(x\) 与 \(y\) （它们属于 \(\mathbf{V} \cap \mathbf{H}_{\beta}\) ），有 \(xy^{-1} \in \mathbf{V}^2 \cap \mathbf{H}_{\beta}\) ；但由 (8)，这蕴含 \(xy^{-1} \in \mathbf{V} \cap \mathbf{H}_{\beta}\) ；换言之，\(\mathbf{V} \cap \mathbf{H}_{\beta}\) 是 \(\mathbf{G}\) 的一个子群，它显然既离散又紧，因而是有限的；但这样一来，由 \(\mathbf{V}_0\) 的取法，这就蕴含确实有 \(\mathbf{V} \cap \mathbf{H}_{\beta} = \{e\}\) 。

我们用反证法，于是假设存在 \(z\) （\(\mathbf{V}^2 -\mathbf{V}\) 的一点）属于 \(\mathbf{H}_{\beta}\) ；由 U 与 V 的取法，在 G 中有 \(\psi (sz^{-1}) + \psi (s)\leqslant \varphi (s)\) ，因为关系式 \(z\notin \mathrm{U}^2\) 蕴含 \(\mathrm{U}z\cap \mathrm{U} = \emptyset\) 。由于

\[
\int \psi (s z ^ {- 1}) d \beta (s) = \int \psi (s) d \beta (s),
\]

由此推出 \(2\beta (\psi)\leqslant \beta (\varphi)\leqslant \alpha (\varphi) + \varepsilon\) ；但我们还有

\[
\beta (\psi) \geqslant \alpha (\psi) - \varepsilon ,
\]

而按构造有 \(\alpha(\varphi) = \alpha(\psi) = \alpha(\{e\})\) 。于是取 \(\varepsilon < \alpha(\{e\})/3\) 便得出矛盾。证毕。

形象地说，满足条件 (L) 的群 G 称为无任意小有限子群的群。*可以证明每个李群都满足条件 (L)；但这一条件并非李群的特征性质；例如，模 \(p\) 同余于 1 的 \(p\) 进整数所成的乘法群就满足 (L)。*

